\documentclass[12pt]{article}

\usepackage[T1]{fontenc}
\usepackage[margin=1in]{geometry}
\usepackage{lmodern}
\usepackage{amsmath,amssymb,amsfonts,bm}
\usepackage[colorlinks=true,linkcolor=blue,citecolor=blue,urlcolor=blue]{hyperref}

\begin{document}

\title{Supplementary Information:\\
Two-dimensional powder diffraction in layered van der Waals films:\\
quantitative X-ray area-detector modeling}
\author{}
\date{}
\maketitle

% Use S-style numbering typical for supplementary material.
\renewcommand{\thesection}{S\arabic{section}}
\renewcommand{\theequation}{S\arabic{equation}}

\section{Incident-beam phase space (finite size and divergence)}
\label{sec:si_beam_phase_space}

The forward model averages kinematic scattering over the finite spatial extent and angular divergence
of the incident beam. We assume a separable Gaussian phase-space density for transverse position
$(x,z)$ (with rms widths $w_x,w_z$) and for small angular deviations $(\delta\theta_x,\delta\theta_z)$
about the nominal beam axis (with rms divergences $\Delta\theta_x,\Delta\theta_z$):
\begin{equation}
  p(x,z,\delta\theta_x,\delta\theta_z)=
  \frac{1}{2\pi w_x w_z}\exp\!\left(-\frac{x^2}{2w_x^2}-\frac{z^2}{2w_z^2}\right)
  \times
  \frac{1}{2\pi \Delta\theta_x \Delta\theta_z}\exp\!\left(-\frac{\delta\theta_x^2}{2\Delta\theta_x^2}-\frac{\delta\theta_z^2}{2\Delta\theta_z^2}\right).
  \label{eq:si_beam_phase_space}
\end{equation}

Detector images are then approximated by Monte Carlo ray sampling: for each ray $n$, we draw
$(x_n,z_n,\delta\theta_{x,n},\delta\theta_{z,n})\sim p$, propagate the ray through the instrument
geometry, and evaluate the kinematic scattering contribution $I_n$ for the corresponding
$(\mathbf{k}_i,\mathbf{k}_f)$. The simulated intensity is the weighted average
\begin{equation}
  I_{\mathrm{det}}\approx \frac{1}{N}\sum_{n=1}^{N} w_n\,I_n,
  \label{eq:si_monte_carlo_average}
\end{equation}
where $w_n$ collects deterministic weights (e.g., absorption/refraction weighting and any optional
per-ray Jacobians) applied after sampling.

\section{Finite footprint and sample-size truncation}
\label{sec:si_footprint_truncation}

At grazing incidence, the vertical beam size produces an elongated illuminated footprint on the
sample surface. A simple scaling for the footprint length is $X_s\sim w_z/\sin\theta_i$, so that
small changes in $\theta_i$ can strongly change the illuminated area.
In the ray model this is handled directly: each sampled ray is intersected with the sample plane,
and rays whose intersection lies outside the finite sample window (a rectangle of in-plane
dimensions $W\times H$) are discarded. Because rays are sampled from the spatial beam profile,
this procedure automatically captures partial illumination, footprint truncation by the sample edge,
and the resulting intensity bias in specular features at very small $\theta_i$.

\section{Refraction and absorption weighting}
\label{sec:si_refraction_absorption}

For rays that intersect the film, we include attenuation through the finite thickness $t$ and
refraction at the interface. Assuming uniform scattering throughout the film and exponential
attenuation along the incident and exit paths, the absorption weight follows from a depth integral:
\begin{equation}
  A(\theta_i,\theta_f)=\int_{0}^{t}\exp\!\left[-\mu z\left(\frac{1}{\sin\theta_i}+\frac{1}{\sin\theta_f}\right)\right]\,dz
  =\frac{1-\exp\!\left[-\mu t\left(\frac{1}{\sin\theta_i}+\frac{1}{\sin\theta_f}\right)\right]}{\mu\left(\frac{1}{\sin\theta_i}+\frac{1}{\sin\theta_f}\right)},
  \label{eq:si_absorption_weight}
\end{equation}
where $\mu$ is the linear absorption coefficient and $\theta_f$ is the grazing-exit angle of the
diffracted wavevector $\mathbf{k}_f$ relative to the sample surface.

Refraction becomes relevant when $\theta_i$ and/or $\theta_f$ are near the critical angle.
We parameterize the optical response by a complex refractive index
$n=1-\delta_{\mathrm{opt}}+i\kappa$, where $\delta_{\mathrm{opt}}$ sets the critical angle and
$\kappa$ sets absorption. In practice, refraction is incorporated by (i) replacing the external
angles $(\theta_i,\theta_f)$ in Eq.~\eqref{eq:si_absorption_weight} by effective internal angles
computed from Snell's law and (ii) multiplying the intensity by Fresnel transmission factors for the
incident and exit waves. This treatment captures the leading GIWAXS refraction/absorption effects
without introducing additional free parameters in the diffraction model.

\section{Wrapped mosaicity distribution}
\label{sec:si_wrapped_distribution}

Orientations are periodic: a tilt angle $\Delta\theta$ is equivalent to $\Delta\theta+2\pi m$ for any
integer $m$. Let $\omega_0(\Delta\theta)$ denote an unwrapped misorientation kernel on the real line
(e.g., the Gaussian--Lorentzian mixture in the main text). The $2\pi$-periodic wrapped distribution is
\begin{equation}
  \omega(\Delta\theta)=\sum_{m=-\infty}^{\infty}\omega_0(\Delta\theta+2\pi m),
  \qquad \Delta\theta\in[-\pi,\pi).
  \label{eq:si_wrapped_definition}
\end{equation}
In computations, the sum is truncated once the neglected terms fall below numerical tolerance.

\section{Bragg-sphere Jacobian for the mosaic surface density}
\label{sec:si_bragg_sphere_jacobian}

For a given reflection, mosaicity redistributes intensity along a Bragg circle in the meridional
$(G_r,G_z)$ plane with line density proportional to $\omega(\vartheta)$, where $\vartheta$ is the
polar angle on the Bragg sphere. Random in-plane rotations about the film normal revolve this line
density around the $G_z$ axis, producing an axially symmetric surface density on the Bragg sphere.

The surface area element of a spherical ring at polar angle $\vartheta$ is
\begin{equation}
  dA = 2\pi R^2 \sin\vartheta\, d\vartheta,
  \label{eq:si_area_element}
\end{equation}
where $R$ is the Bragg-sphere radius. If the intensity assigned to the corresponding polar interval
is $dI = I_0\,\omega(\vartheta)\,d\vartheta$, then the surface density is
\begin{equation}
  \sigma(\vartheta)=\frac{dI}{dA}=\frac{I_0}{2\pi R^2\sin\vartheta}\,\omega(\vartheta),
  \label{eq:si_sigma_jacobian}
\end{equation}
which yields the $1/\sin\vartheta$ Jacobian factor in the main-text normalization.

\end{document}
